\documentclass[11pt]{article}
\usepackage[utf8]{inputenc}
\usepackage[T1]{fontenc}
\usepackage{amsmath,amssymb,amsthm}
\usepackage{geometry}
\usepackage{booktabs}
\usepackage{hyperref}
\geometry{margin=2.5cm}

\newtheorem{theorem}{Theorem}[section]
\newtheorem{proposition}[theorem]{Proposition}
\newtheorem{corollary}[theorem]{Corollary}
\theoremstyle{remark}
\newtheorem{remark}[theorem]{Remark}

\newcommand{\lean}[1]{\texttt{#1}}
\newcommand{\curl}{\operatorname{curl}}

\title{A Machine-Checked Audit of the FB(S$^3$)R Archive in Lean~4:\\
Killing--Beltrami Algebra, $\zeta$-Constants, Bell--Hopf Rigidity,\\
Zeckendorf--Collatz Bounds, and Six Corrected Claims}
\author{FB(S$^3$)R Research Archive}
\date{26 November 2026}

\begin{document}
\maketitle

\begin{abstract}
We report a formal verification, in the Lean~4 proof assistant with the Mathlib library,
of a selection of mathematical statements from the FB(S$^3$)R preprint archive. Twelve Lean
modules contain seventy-five theorems. Grouped by content and leaving out trivial auxiliary
lemmas, they yield 47 results: 6 corrected claims and 41 verified statements. None uses \lean{sorry}, \lean{admit},
\lean{native\_decide} or any user-declared axiom. Every theorem depends only on Lean's
standard axioms \lean{propext}, \lean{Classical.choice} and \lean{Quot.sound}.
The verified material comprises:
\begin{itemize}
\item the algebraic core of the Killing--Beltrami theory of Navier--Stokes flow on the round
three-sphere: the representation-defect theorem, the polarization coefficients, the curl-shell
spectral identities, and the energy law of Killing precession;
\item the Lamb--Beltrami mechanism for differentiable fields, and explicit Beltrami
(ABC) steady Euler solutions;
\item the $\zeta$-function shell sum and the exact residue $1/(3\pi^2)$ of the supercritical
embedding constant;
\item the CHSH bound for every local hidden-variable model, the Tsirelson bound for the singlet
kernel, and Chern-charge rigidity $q=\pm1$ for the Hopf holonomy representation;
\item the Zeckendorf height bounds for the Collatz map;
\item the causal-spectral identities of the round carrier: geodesic-ball volume, the causal
fraction, and the scalar, transverse-traceless and Dirac capacities;
\item Berger-sphere curvature and the Ricci-flow ratio equation, the GHZ branch-balancing
factorization, the dilation algebra of valuation profiles, and the vacuum-shift kernel of
trace-free gravity.
\end{itemize}
We also give machine-checked counterexamples to six claims made in the archive.
We state explicitly which statements were not formalized, and why.
No open problem (Collatz, Navier--Stokes regularity, the Yang--Mills mass gap, local
connectivity of the Mandelbrot set) is claimed to be solved.
\end{abstract}

\section{Method and trust base}
The development is the Lake project \lean{Lean/FBS3R-Lean} (Lean~4.36.0, Mathlib
\lean{v4.36.0}). The command \lean{lake build FBS3R} compiles it, and
\lean{lake env lean scripts/Axioms.lean} prints the axiom dependencies of every
principal theorem. All of them report \lean{[propext, Classical.choice, Quot.sound]}.

Mathlib does not yet provide Hodge theory or a curl operator on Riemannian three-manifolds.
Statements about the round $S^3(R)$ are therefore split into two parts:
\begin{enumerate}
\item the \emph{geometric input}, which is assumed as a hypothesis: the signed curl shells
$E_{k,\sigma}$ have eigenvalue $\mu_{k,\sigma}=\sigma(k+2)/R$, the Killing algebra
$\mathfrak{so}(4)=\mathfrak{su}(2)_+\oplus\mathfrak{su}(2)_-$ acts on them with spins
$j_\tau=(k+1+\tau\sigma)/2$, and the relevant invariant kernels are constant;
\item the \emph{algebra and analysis} derived from that input, which is proved in full.
\end{enumerate}
The geometric input consists of classical facts. Its formalization is left to future work.

\section{Killing--Beltrami algebra}
\begin{theorem}[Representation defect; \lean{KillingBeltrami.defect}]
Let $P:L_1\to\mathrm{End}(M)$ and $Q:L_2\to\mathrm{End}(M)$ be Lie-algebra homomorphisms
with $[P x,Q y]=0$ for all $x,y$. For scalars $a,b$ put $D(x,y)=aP(x)+bQ(y)$. Then
\[
[D(m),D(n)]-D([m_1,n_1],[m_2,n_2])=(a^2-a)P[m_1,n_1]+(b^2-b)Q[m_2,n_2].
\]
For $a=k/(k+2)$ and $b=(k+4)/(k+2)$ the two coefficients are $-2k/(k+2)^2$ and
$2(k+4)/(k+2)^2$. Both are nonzero for $k\ge1$ (\lean{defect\_coeffs\_ne\_zero}).
\end{theorem}
This is the archive's identity $\Omega_{k,\sigma}$, proved without the $S^3$ realization.
The formula $D=\mathrm{ad}+\tfrac{2}{k+2}(\mathrm{ad}_{m_{-\sigma}}-\mathrm{ad}_{m_\sigma})$
is \lean{D\_minus\_ad}.

\begin{proposition}[Coefficients and spectra]
For $\sigma,\tau\in\{\pm1\}$, $R\neq0$ and $k\in\mathbb N$:
\begin{enumerate}
\item $(\mu_{k,\sigma}-\mu_{0,\tau})/(2\mu_{k,\sigma})=(k+2-2\tau\sigma)/(2(k+2))$
(\lean{KB\_coefficient}), with the chiral values $k/(k+2)$ and $(k+4)/(k+2)$;
\item $(2j_+ + 1)(2j_- + 1)=(k+1)(k+3)$, and
$\tfrac{4}{R^2}\sum_\tau j_\tau(j_\tau+1)=2\mu_{k,\sigma}^2$, that is,
$\Delta_{\mathrm{iso}}=2\curl^2$;
\item $\mu_{k,\sigma}^2-4/R^2=k(k+4)/R^2$, which is the Stokes eigenvalue;
\item the closed structure coefficients equal $k^2(k+4)/(4(k+2))$ and $k(k+4)^2/(4(k+2))$, and
$(k+1)(k+3)/(k+2)^2=1-(k+2)^{-2}$.
\end{enumerate}
\end{proposition}

\begin{theorem}[Killing precession does no work; \lean{Precession}]
Let $E$ be a real inner-product space and suppose $v'=-cv-D(t)v+F(t)$, where each $D(t)$ is
skew. Then $\tfrac{d}{dt}\tfrac12\|v\|^2=-c\|v\|^2+\langle F,v\rangle$. If moreover $F=0$, then
$\|v(t)\|=e^{-ct}\|v(0)\|$ for all $t$.
\end{theorem}

\section{Lamb identity, Beltrami fields and ABC flows}
\begin{theorem}[\lean{Lamb.lamb}, \lean{Lamb.beltrami\_advection\_is\_gradient}]
Let $u:\mathbb R^3\to\mathbb R^3$ be Fr\'echet differentiable at $x$. Then
$(u\cdot\nabla)u=\nabla\tfrac12|u|^2-u\times\curl u$ at $x$. In particular, if
$\curl u=\lambda u$ at $x$, the advection term at $x$ is the gradient of $\tfrac12|u|^2$.
The polarized identity $\nabla_uv+\nabla_vu=\nabla\langle u,v\rangle-u\times\curl v
-v\times\curl u$ also holds at the level of Jacobians.
\end{theorem}
\begin{theorem}[ABC flows; \lean{ABC}]
For all $A,B,C\in\mathbb R$, the field
$U=(A\sin z+C\cos y,\ B\sin x+A\cos z,\ C\sin y+B\cos x)$ satisfies $\curl U=U$ and
$\operatorname{div}U=0$. It also satisfies $(U\cdot\nabla)U=\nabla\tfrac12|U|^2$, so it is a
steady Euler solution with pressure $p=-\tfrac12|U|^2$.
\end{theorem}

\section{$\zeta$-constants of the shell barriers}
Let $Z(s)=\sum_{n\ge1}n^{-s}$. Lean proves $Z(s)=\zeta(s)$ for real $s>1$
(\lean{Z\_eq\_riemannZeta}).
\begin{theorem}[\lean{ZetaShell}]
\begin{enumerate}
\item For $s>3$: $\sum_{k\ge0}2(k+1)(k+3)(k+2)^{-s}=2(\zeta(s-2)-\zeta(s))$.
\item With $V_R=2\pi^2R^3$: $2R^s/(3V_R)=R^{s-3}/(3\pi^2)$.
\item For $C_s(R)=\frac{R^{s-3}}{3\pi^2}(\zeta(s-2)-\zeta(s))$:
$\lim_{s\downarrow3}(s-3)C_s(R)=1/(3\pi^2)$ for every $R>0$.
\end{enumerate}
\end{theorem}

\section{Projector bounds, entropy and other checks}
\begin{theorem}[\lean{Projector}, \lean{Entropy}]
Let $(e_i)_{i\le N}$ be orthonormal in $L^2(\mu)$, where $\mu(X)=V$, and suppose
$\sum_ie_i(x)^2\equiv K$. Then:
\begin{itemize}
\item $K=N/V$;
\item $w(x)^2\le\|a\|^2N/V$ for $w=\sum_i a_ie_i$, with equality attained;
\item $\|a\|^2/w(x)^2\ge V/N$.
\end{itemize}
In addition, the conditional entropy bound $\mathcal S\ge\delta\log(M/C)$ holds for every
finite probability vector satisfying the hypotheses of the archive.
\end{theorem}

\section{Bell correlations and the Hopf bundle}
\begin{theorem}[\lean{Bell}]
\begin{enumerate}
\item For every probability space and all measurable $\pm1$-valued responses,
$|E(a,b)+E(a,b')+E(a',b)-E(a',b')|\le2$.
\item For unit vectors in any real inner-product space, the kernel
$E(a,b)=-\langle a,b\rangle$ gives $|S|\le2\sqrt2$. The PRA settings give exactly
$S=-2\sqrt2$. Hence no local hidden-variable model reproduces this kernel.
\item (Chern rigidity) For $q\in\mathbb Z$: $\cos(q(\pi-\theta))=-\cos\theta$ holds for all
$\theta\in[0,\pi]$ if and only if $q=\pm1$.
\item The Hopf map $(z,w)\mapsto(2z\bar w,|z|^2-|w|^2)$ sends $S^3$ to $S^2$ and is
constant on $U(1)$ fibres.
\end{enumerate}
\end{theorem}

\section{Zeckendorf heights and the Collatz map}
Let $h(n)$ be the largest index $k$ with $F_k\le n$. This is Mathlib's \lean{greatestFib}.
\begin{theorem}[\lean{Collatz}]
\begin{enumerate}
\item $h(\lfloor n/2\rfloor)+1\le h(n)$ for $n\ge1$.
\item $h(3n+1)\le h(n)+3$.
\item $h((3n+1)/2^m)+m\le h(n)+3$ whenever $2^m\mid 3n+1$.
\end{enumerate}
The worked examples of the naive-Lyapunov and coarse-certificate no-go results are verified
as well. The Fibonacci intervals $[F_n,F_{n+1})$, $n\ge2$, partition $\mathbb Z_{\ge1}$, and
$F_{n+1}/F_n\to\varphi$.
\end{theorem}
These are local bounds only. They do not imply the Collatz conjecture, and none of the
preprints claims that they do.

\section{Causal--spectral carrier identities}
\begin{theorem}[\lean{CausalSpectral}]
On the round sphere $S^3(R)$ the following hold.
\begin{enumerate}
\item $\int_0^s4\pi R^2\sin^2(\rho/R)\,d\rho=2\pi R^3\bigl(s/R-\tfrac12\sin(2s/R)\bigr)$,
which equals $2\pi^2R^3$ at $s=\pi R$.
\item The causal fraction $\Gamma(\tau)=(\tau-\tfrac12\sin2\tau)/\pi$ satisfies
$\Gamma'=\tfrac2\pi\sin^2\tau$, $\Gamma(\tau)+\Gamma(\pi-\tau)=1$ and $\Gamma(\pi/2)=\tfrac12$.
\item The scalar cutoff satisfies $n(n+2)\le x^2\iff n+1\le\sqrt{1+x^2}$, and
$\sum_{n<M}(n+1)^2=M(M+1)(2M+1)/6$.
\item $\sum_{k=2}^M2(k-1)(k+3)=M(M-1)(2M+11)/3$.
\item $2\sum_{n=0}^M(n+1)(n+2)=2(M+1)(M+2)(M+3)/3$.
\end{enumerate}
\end{theorem}

\section{Further verified statements}
\begin{theorem}[\lean{Checks}]
\begin{enumerate}
\item (Hierarchical Bethe graphs.) In the maximally coupled block adjacency
$\begin{pmatrix}A&\mathbf 1\mathbf 1^T\\\mathbf 1\mathbf 1^T&0\end{pmatrix}$ with a new block of
size $k$, the kernel has dimension at least $k-1$.
\item (Berger spheres.) With $x=a^2$ and $z=c^2$:
\begin{itemize}
\item $R=8/x-2z/x^2$ and $|\mathrm{Ric}|^2=4[2(2x-z)^2+z^2]/x^4$;
\item $4|\mathrm{Ric}|^2-R^2=4(16x^2-24xz+11z^2)/x^4$, which tends to $64/a^4$ as $c\to0$;
\item $\dot x=-2r_1x$ and $\dot z=-2r_3z$ imply $\dot u=8u(1-u)/x$ for $u=z/x$.
\end{itemize}
\item (GHZ balancing.) $\frac{2cs|\Gamma|}{c^2A+s^2B}=q\frac{2r}{1+r^2}=\frac{q}{\cosh\ln r}$,
and $\frac{2r}{1+r^2}\le1$ with equality if and only if $r=1$.
\item (Dilation algebra.)
\begin{itemize}
\item From $E=\varepsilon V$, $V=2\pi^2A^3$ and $dE=-p\,dV$ it follows that
$p=-\varepsilon-\tfrac A3\varepsilon'$.
\item $A\frac{d}{dA}A^n=nA^n$.
\item The integer roots of $(z-1)z(z+1)(z+2)(z+3)$ are exactly $\{1,0,-1,-2,-3\}$.
\item These are the exponents of the round valuation density.
\end{itemize}
\item (Trace-free gravity.) For every invertible $4\times4$ metric $g$:
$P_{\mathrm{tr}}(T-\rho_0g)=P_{\mathrm{tr}}(T)$.
\end{enumerate}
\end{theorem}
\begin{remark}[Numerical check]
The value $q_0=\pi\varphi^{584}/(1+\pi/50)=3.30725146071398\ldots\times10^{122}$ is reproduced.
This checks the arithmetic only. It does not derive $\Lambda$, because both the shell number
$292$ and the factor $1+\pi/50$ are inputs of the paper.
\end{remark}

\section{Corrected claims}
\begin{proposition}[Refined odd step is false as stated; \lean{refined\_odd\_step\_counterexample}]
Take $n=11$. Then $6\le\log_\varphi(11\sqrt5)<6+\tfrac23$, so $11$ is not ``top-third''.
Yet $h(34)-h(11)=9-6=3$. An exhaustive search over odd $n\le2\cdot10^6$ finds no other
violation. The statement can be repaired by adding the hypothesis ``for $n$ sufficiently
large''. The drift inequality $k_{J+1}-k_J\le(2+I_J)-A_J$ then holds only eventually.
\end{proposition}

\begin{proposition}[Uniqueness of the Fibonacci stratification is vacuous; \lean{FibStrat}]
Let $(a_n)$ be strictly increasing with $a_0=1$ and assume that each sum $a_m+a_n$ is again
some $a_k$. Then $a_n=n+1$, so $a_{n+1}/a_n\to1$. Consequently no partition satisfies the
hypotheses of the uniqueness theorems in the preprints of 15 and 19 April 2026. Moreover,
the Fibonacci numbers themselves are not additively closed: $F_2+F_5=6$.
\end{proposition}

\begin{proposition}[Bell state versus singlet; \lean{bell\_Phi\_plus}]
For $|\Phi^+\rangle=(|00\rangle+|11\rangle)/\sqrt2$:
$\langle\sigma_x\otimes\sigma_x\rangle=+1$, $\langle\sigma_y\otimes\sigma_y\rangle=-1$ and
$\langle\sigma_z\otimes\sigma_z\rangle=+1$. So $E(a,b)=-\cos\theta_{ab}$, as claimed in the
preprint of 15 April 2026, is false for $\Phi^+$. It holds for the singlet $\Psi^-$, where all
three values equal $-1$; the PRA version uses the correct state.
\end{proposition}

\begin{proposition}[Yang--Mills threshold; \lean{no\_three\_in\_curl\_or\_stokes\_spectrum}]
The value $3/R^2$ of the preprint of 19 April 2026 lies neither in the curl-squared spectrum
$(k+2)^2/R^2$, whose minimum is $4/R^2$, nor in the Stokes spectrum $k(k+4)/R^2$. Both
spectra are used elsewhere in the same archive. Either the operator must be specified, or the
threshold must be corrected.
\end{proposition}

\begin{proposition}[Binet formula of the Bethe-graph paper; \lean{paper\_binet\_wrong}]
For $F_1=3$, $F_2=5$, $F_{n+1}=F_n+F_{n-1}$, the formula
$F_n=(\varphi^{n+2}-\psi^{n+2})/\sqrt5$ gives $F_1=2$. The correct statement is
$F_n=\mathrm{Fib}(n+3)=(\varphi^{n+3}-\psi^{n+3})/\sqrt5$ (\lean{F\_binet}). The ratio limit
$F_{n-1}/F_n\to\varphi^{-1}$ is unaffected.
\end{proposition}

\begin{remark}[Mandelbrot local connectivity]
The preprint of 5 March 2025 does not prove the MLC conjecture, which remains open. Two
steps fail.
\begin{itemize}
\item Trivial fundamental group does not imply local connectivity: the comb space is
contractible but not locally connected.
\item A subset of a connected disc need not be connected. Lean checks the counterexample
$\{0,1\}\subset[0,1]$ (\lean{subset\_of\_connected\_not\_connected}).
\end{itemize}
\end{remark}

\section{Not formalized}
\begin{itemize}
\item The spectral geometry of $S^3$: curl eigenspaces, Leray projection, and the
classification of reducing spaces.
\item The vector Laplacian $\Delta U=-U$ for ABC flows, which is needed for the decaying
Navier--Stokes solution $e^{-\nu t}U$.
\item The conditional Euclidean blow-up transfer, which depends on five external hypotheses.
\item Cosmological, empirical and ontological claims, which are not mathematical theorems.
\end{itemize}

\section{Ledger}
The file \lean{RESULTS.md} in the Lean project lists all 47 results with their sources and
Lean names: A1--A6 are corrections and B1--B41 are verifications. Trivial auxiliary lemmas are
listed in the code but not counted. No result in this ledger settles an open problem.

\end{document}
